\section{Lời giải thích như một hàm xấp xỉ}
\label{sec:ch13-ham-xap-xi}

\index{lời giải thích!mô hình hóa thành hàm xấp xỉ|(}%
Nếu với một kết quả dạng chặn thì chỗ cần soi là các định nghĩa chứ không phải
phần chứng minh, thì phải có đủ bộ định nghĩa trước mắt đã. Mục này dựng chúng,
và cả cái trần lẫn mọi giới hạn của nó đều nằm sẵn trong đây.

Bài báo bắt đầu bằng hai định nghĩa ngắn. Mô hình cần giải thích là một hàm
$f : \mathcal{X} \to \mathcal{Y}$ từ không gian đầu vào sang không gian đầu ra.
Lời giải thích cho $f$ là một hàm $g : \mathcal{X} \to \mathcal{Y}$ xấp xỉ $f$
và được coi là dễ diễn giải với con người theo một tiêu chí nào
đó~\autocite{p26limits}. Hai định nghĩa ấy đặt ra hai đòi hỏi và bỏ ngỏ đúng một
chỗ: cái gì làm cho một hàm là dễ diễn giải. Toàn bộ phần còn lại của bài là một
câu trả lời cho chỗ bỏ ngỏ ấy, và câu trả lời là một hạn mức tính bằng bit.

Nhưng trước khi tới hạn mức thì phải nói vì sao cần một thước đo độ phức tạp
kiểu này, chứ không phải một thước đo hiển nhiên hơn. Cách hiển nhiên là đếm
tham số: một mô hình tuyến tính mười hệ số thì đơn giản hơn một mạng nơ-ron
triệu tham số. Phép đếm ấy hỏng ở chỗ nó phụ thuộc vào cách viết. Một ma trận
một triệu phần tử mà mọi phần tử đều bằng nhau thì viết ra hết là một triệu số,
mà mô tả bằng lời chỉ là một câu: lấy một số rồi lặp lại một triệu lần. Đếm tham
số cho hai kết quả rất khác nhau cho cùng một đối tượng, tùy người đọc nhìn vào
bản viết nào. Cái ta muốn đo là lượng thông tin thật sự cần để chỉ định hàm ấy,
không phụ thuộc cách trình bày, và đó chính là đại lượng mà AIT đo.

\index{độ phức tạp Kolmogorov}%
Thước đo ấy là \tn{độ phức tạp Kolmogorov}{Kolmogorov complexity}, và định
nghĩa của nó là định nghĩa 2.3 của bài~\autocite{p26limits}.

\begin{definition}[Độ phức tạp Kolmogorov]
\label{def:ch13-do-phuc-tap-kolmogorov}
Cho một \tn{máy Turing phổ dụng}{universal Turing machine} $U$. Độ phức tạp
Kolmogorov $K(g)$ của một hàm $g$ là độ dài, tính bằng bit, của chương trình
ngắn nhất mà $U$ chạy được để tính $g$:
\[
  K(g) \;=\; \min_{p \in \{0,1\}^{*}} \{\, |p| \;:\; U(p,x) = g(x) \text{ với mọi } x \in \mathcal{X} \,\}.
\]
\end{definition}

Định nghĩa~\ref{def:ch13-do-phuc-tap-kolmogorov} xử lý đúng ví dụ ma trận ở
trên: chương trình
ngắn nhất sinh ra ma trận toàn số giống nhau ngắn hơn hẳn chương trình liệt kê
một triệu số, nên độ phức tạp Kolmogorov của nó nhỏ, bất kể người ta viết nó ra
dưới dạng nào. Với mục đích của bài, đại lượng này được đọc là lượng tài nguyên
nhận thức tối thiểu cần để nắm được hành vi của hàm.

Một dấu hỏi hiện ra ngay: $K$ phụ thuộc vào máy $U$ đã chọn, mà có vô số máy phổ
dụng. Bài dẫn định lý bất biến, một kết quả kinh điển của lý thuyết này, để trả
lời: với hai máy phổ dụng bất kỳ $U_1$ và $U_2$, tồn tại một hằng số
$c_{U_1,U_2}$ không phụ thuộc $g$ sao cho hai giá trị $K$ chênh nhau không quá
hằng số ấy~\autocite{p26limits}. Chứng minh chỉ là viết một chương trình thông
dịch: máy $U_2$ mô phỏng được $U_1$ bằng một đoạn mã cố định, nên bất kỳ chương
trình nào cho $U_1$ cũng thành một chương trình cho $U_2$ dài hơn đúng độ dài
đoạn mã ấy. Hằng số không đổi trong khi $K(g)$ tăng theo độ phức tạp của $g$,
nên với những hàm đủ phức tạp thì việc chọn máy nào không còn quan trọng. Chú ý
rằng lập luận này chỉ đúng theo nghĩa tiệm cận: với một hàm cụ thể và một hằng
số đủ lớn, $K$ vẫn là một đại lượng xác định sai khác một lượng không kiểm soát
được.

Cùng lúc đó, $K$ không tính được. Không có thuật toán nào nhận vào một hàm và
trả ra độ phức tạp Kolmogorov của nó, và đây là một trong những kết quả cổ điển
nhất của lý thuyết. Bài báo ghi nhận điều này ở trang 6 và trả lời rằng $K$ ở
đây đóng vai trò nền lý thuyết, còn với các lớp mô hình cụ thể thì nó xấp xỉ
được, theo cách mà mục~\ref{sec:ch13-tran-cao-bao-nhieu} sẽ
đọc~\autocite{p26limits}. Hệ quả của tình trạng ấy đối với việc dùng các chặn
này trong thực hành là một chuyện riêng, và
mục~\ref{sec:ch13-dieu-chua-giai-quyet} nói tới.

\index{lớp khả diễn giải}%
Với thước đo trong tay, hạn mức viết ra rất gọn. Với một ngưỡng
$b \in \mathbb{N}$, đặt
\[
  \mathcal{G}_b \;=\; \{\, g : \mathcal{X} \to \mathcal{Y} \;\mid\; K(g) \le b \,\},
\]
tức tập mọi hàm mô tả được trong $b$ bit. Bài gọi đây là \tn{lớp khả diễn
giải}{interpretability class} và đọc $b$ là dung lượng nhận thức mà người đọc có
sẵn~\autocite{p26limits}. Đây là bước mô hình hóa quan trọng nhất trong cả bài,
vì nó biến một khái niệm thuộc về con người thành một con số, và mọi kết quả về
sau đều thừa hưởng cả sức mạnh lẫn điểm yếu của bước ấy.

Một quan sát đếm được về lớp này sẽ dùng lại nhiều lần. Số chương trình dài
không quá $b$ bit là $\sum_{i=0}^{b} 2^i = 2^{b+1} - 1$, và hai chương trình khác
nhau có thể tính cùng một hàm nhưng một chương trình không tính được hai hàm,
nên $|\mathcal{G}_b| \le 2^{b+1} - 1$~\autocite{p26limits}. Số hàm dễ đọc là hữu
hạn và tăng theo hàm mũ của hạn mức, còn số hàm mà một mô hình có thể là thì lớn
hơn thế rất nhiều. Mọi chặn định lượng của
mục~\ref{sec:ch13-tran-cao-bao-nhieu} đều là phép đếm này viết cho tử tế; cái
trần ở mục~\ref{sec:ch13-tran-va-chung-minh} thì đến từ một chỗ khác và mục ấy
nói rõ chỗ nào.

\index{sai số xấp xỉ}%
Còn thiếu vế thứ hai: xấp xỉ dở tới mức nào. Bài đặt $\rho(y_1,y_2)$ là khoảng
cách giữa hai đầu ra và định nghĩa hai cách đo sai số xấp xỉ, tức khoảng cách
giữa lời giải thích và mô hình xét trên đầu ra~\autocite{p26limits}. Sai số kỳ
vọng
$\mathcal{E}_{\mathcal{D}}(f,g) = \mathbb{E}_{x \sim \mathcal{D}}[\rho(f(x),g(x))]$
lấy trung bình chỗ lệch trên phân phối đầu vào $\mathcal{D}$, tức phân phối các
đầu vào gặp trong thực tế. Sai số trường hợp xấu nhất
$\mathcal{E}_{\infty}(f,g) = \sup_{x \in \mathcal{X}} \rho(f(x),g(x))$ lấy chỗ
lệch lớn nhất trên toàn không gian. Vì trung bình không vượt quá cực đại nên
$\mathcal{E}_{\mathcal{D}} \le \mathcal{E}_{\infty}$ với mọi $\mathcal{D}$: một
bảo đảm về trường hợp xấu nhất kéo theo bảo đảm về trung bình, còn chiều ngược
lại thì không. Phân biệt này không phải chi tiết kỹ thuật. Một lời giải thích
đúng trên hầu hết đầu vào và sai hẳn trên một vùng nhỏ có sai số kỳ vọng nhỏ và
sai số trường hợp xấu nhất lớn, mà vùng nhỏ ấy đúng là chỗ
chương~\ref{ch:giai-thich-duoi-tan-cong} sẽ tìm tới.

Hai đại lượng còn lại là hai cách nhìn cùng một sự đánh đổi, và bài đặt tên cho
cả hai: \tn{hàm sai số giải thích}{explanation error function} và \tn{hàm độ
phức tạp giải thích}{explanation complexity function}.

\begin{definition}[Hàm sai số giải thích]
\label{def:ch13-ham-sai-so}
Với mô hình $f$, ngưỡng $b$ và một trong hai cách đo sai số $\mathcal{E}$,
\[
  \varepsilon_f(b) \;=\; \inf_{g \in \mathcal{G}_b} \mathcal{E}(f,g)
\]
là sai số nhỏ nhất mà một lời giải thích không quá $b$ bit đạt được.
\end{definition}

\begin{definition}[Hàm độ phức tạp giải thích]
\label{def:ch13-ham-do-phuc-tap}
Với mô hình $f$ và ngưỡng sai số $\delta > 0$,
\[
  \kappa_f(\delta) \;=\; \min \{\, b \in \mathbb{N} \;\mid\; \exists\, g \in \mathcal{G}_b : \mathcal{E}(f,g) \le \delta \,\}
\]
là số bit ít nhất đủ để đạt sai số không quá $\delta$.
\end{definition}

Hai định nghĩa này là định nghĩa 2.12 và 2.16 của bài~\autocite{p26limits}. Cái
thứ nhất cố định ngân sách rồi hỏi sai số; cái thứ hai cố định sai số rồi hỏi
ngân sách. Bài chứng minh chúng là hai nửa của một phát biểu: $\varepsilon_f(b)
\le \delta$ khi và chỉ khi $b \ge \kappa_f(\delta)$~\autocite{p26limits}. Cả hai
đều đơn điệu theo chiều dễ đoán, nới ngân sách thì sai số không tăng và nới sai
số thì ngân sách không tăng, nên đồ thị của $\varepsilon_f$ là một đường không
tăng và mọi thứ chương này nói đều đọc được trên đường ấy.

\begin{figure}[htbp]
  \centering
  \begin{tikzpicture}[
  every node/.style={font=\footnotesize, align=center},
  axis/.style={-{Stealth[length=2mm]}, thick, black!70},
  curve/.style={very thick, black!75},
  cut/.style={thick, black!55, densely dashed},
  point/.style={thick, black!55},
  empty/.style={draw=black!45, dotted, thick, fill=black!2},
  lbl/.style={font=\scriptsize, align=center, black!55},
  tick/.style={font=\scriptsize, black!70}
]
  % the region above the curve: every (budget, error) pair an explanation can reach
  \fill[black!7]
    (0.55,4.2) .. controls (1.6,1.35) and (3.0,0.28) .. (6.2,0)
    -- (7.1,0) -- (7.1,4.5) -- (0.55,4.5) -- cycle;

  \draw[axis] (0,0) -- (7.6,0) node[tick, below left=1mm and 0mm] {$b$};
  \draw[axis] (0,0) -- (0,4.9) node[tick, above left=0mm and 1mm] {$\delta$};

  % the box no explanation reaches: budget under the human bound, error under the threshold
  \draw[empty] (0,0) rectangle (2.1,0.62);

  \draw[curve] (0.55,4.2) .. controls (1.6,1.35) and (3.0,0.28) .. (6.2,0);
  \node[lbl, anchor=west, text width=20mm] at (3.5,1.15) {$\varepsilon_f(b)$};

  \draw[cut] (2.1,0) -- (2.1,4.5);
  \draw[cut] (0,0.62) -- (7.1,0.62);

  \filldraw[black!75] (6.2,0) circle (1.5pt);

  \node[tick, below=1.5mm] at (2.1,0) {$b_{\text{người}}$};
  \node[tick, below=1.5mm] at (6.2,0) {$K(f)$};

  \node[lbl, anchor=east, text width=26mm] at (7.6,3.3)
    {miền khả thi: mọi cặp\\lấy được bằng một\\lời giải thích};
  \node[lbl, anchor=west, text width=30mm] at (2.45,-1.05)
    {ô rỗng: không lời giải thích nào\\vừa dưới $b_{\text{người}}$ bit\\vừa dưới sai số $\delta$};
  \draw[point, -{Stealth[length=1.6mm]}] (2.6,-0.75) -- (1.2,-0.05);
\end{tikzpicture}

  \caption{Hai nét đứt là hai lựa chọn tự do, theo quy ước mà
  hình~\ref{fig:ch08-vong-lap-do} đặt: người dùng chọn cả hạn mức bit lẫn ngưỡng
  sai số. Nét chấm là thứ đang thiếu, theo quy ước mà
  hình~\ref{fig:ch10-vong-bien-the} đặt: ô rỗng là chỗ mà không lời giải thích
  nào của $f$ rơi vào được. Điểm $K(f)$ trên trục ngang là chỗ
  đường cong chạm không, tức chỉ có một chương trình tính đúng $f$ mới giải
  thích nó không sai chút nào. Bản vẽ của sách từ các định nghĩa 2.12, 2.16 và
  4.2 của bài báo~\autocite{p26limits}, vốn không có hình nào.}
  \label{fig:ch13-mien-kha-thi}
\end{figure}

Hình~\ref{fig:ch13-mien-kha-thi} vẽ nó. Trục ngang là ngân sách bit, trục dọc là
sai số, đường cong là $\varepsilon_f$, và vùng phía trên đường cong là tập các
cặp mà một lời giải thích thật sự lấy được: bài gọi đó là \tn{miền khả
thi}{feasibility region} khi bàn tới quy định, và
mục~\ref{sec:ch13-tran-noi-gi} dùng lại tên ấy. Hai đường đứt cắt hình là hai
ràng buộc mà một người dùng có thể đặt ra: một hạn mức bit ứng với sức đọc của
con người, và một ngưỡng sai số coi là chấp nhận được. Ô chữ nhật mà hai đường
ấy cắt ra ở góc dưới bên trái nằm hoàn toàn dưới đường cong, nghĩa là rỗng.
Mục~\ref{sec:ch13-tran-va-chung-minh} chứng minh rằng với những mô hình đủ phức
tạp thì ô ấy rỗng thật.

\index{lời giải thích!mô hình hóa thành hàm xấp xỉ|)}%
